\chapter{Conclusions}
\label{chap:conclusions}

Reinforcement learning offers a means of acquiring complex behaviour through interaction rather than specifying that behaviour directly. Its practical use, however, often depends on substantial specification elsewhere in the learning system. Before an agent can be trained, a sufficiently accurate and interactable simulator must generally be constructed, and that simulator must expose the agent to a sufficiently diverse set of scenarios. This thesis has referred to these requirements collectively as the \emph{simulation bottleneck}. The bottleneck limits the domains in which RL can be applied economically and creates a persistent risk that agents succeed within their engineered training conditions while failing to generalise beyond them.

This thesis has investigated the simulation bottleneck at three levels. At the level of the agent, it studied how the representations learned by actor--critic algorithms may overfit to the environment instances encountered during training when the number of available instances is limited. At the level of the training distribution, it examined how the available set of simulator instances can be expanded through generative modelling without losing correspondence to the target environment distribution. At the level of the simulator, it explored how action-conditioned visual experience can itself be generated by a learned model that maintains a persistent representation of the world. These levels form a progression from improving learning within a fixed simulator, to reducing the manual work required to populate that simulator, and finally towards learning the simulator itself from data.

A central conclusion of the thesis is that generalisation in RL is not solely a property of the policy or learning algorithm. It emerges from the interaction between the information retained by the agent, the distribution of environments on which it is trained, and the mechanism that produces its experience. An architecture that promotes transferable representations may still fail if the available training instances are too few or insufficiently representative of the target distribution.

The findings also show that generative modelling can intervene at different levels of the simulation bottleneck. DRED targets its second condition by transforming a finite dataset of simulator parameters into a broader distribution of training scenarios grounded in the target distribution, thereby reducing the need to specify environment variations manually. PERSIST targets its first condition by learning an interactable approximation of the environment from recorded trajectories, reducing dependence on a hand-engineered simulator. These interventions are complementary rather than instances of the same mechanism: one expands the scenarios available within an existing simulator, while the other moves towards replacing the simulator itself. Together, they demonstrate how generative modelling can progressively relax the simulation bottleneck, from reducing manual scenario design to reducing manual simulator construction.

The work presented here does not eliminate the simulation bottleneck. DRED still requires an existing simulator in which generated instances can be executed, while PERSIST does not yet establish that policies trained entirely within its generated experience will transfer reliably to the underlying environment. Instead, the thesis develops a hierarchy of concrete interventions and clarifies the requirements for further progress. At the level of the agent, improving generalisation requires representations that retain instance-agnostic information. At the level of the training distribution, environment generation can encourage such representations by exposing the agent to a broader and more diverse set of instances, provided that generation remains grounded in the target distribution. At the level of the simulator, learned world models offer a route towards training agents without hand-engineered simulation and may help bridge the sim-to-real gap, but they must preserve a coherent world state over the extended horizons required for control.


\section{Main Findings}
\label{sec:conclusion_main_findings}

In the following, we summarise the central findings of this work.

\paragraph{Chapter~2: Actor and critic representations specialise, and their coupling affects generalisation.}
Chapter~2 investigated how learned representations mediate instance overfitting in actor--critic reinforcement learning. It distinguished generalisation from training performance: an agent may perform well on the finite environment instances encountered during optimisation without acquiring behaviour that transfers to unseen instances drawn from the same target distribution. To characterise this failure, the chapter introduced mutual-information-based metrics that quantify how much a learned representation encodes about inverse dynamics, forward dynamics, state value, and training-instance identity. The chapter further established that this last quantity defines an upper bound on the generalisation error, thereby providing a formal connection between instance-specific representations and zero-shot transfer.

The empirical results show that actor and critic representations systematically specialise when learned separately. Actor representations preferentially retain information that supports action selection, whereas critic representations encode more information about value and future states and, consequently, tend to extract more instance-specific information. When the actor and critic share a representation, features that help the critic distinguish between individual training environments are also made available to the actor, encouraging it to rely on instance-specific features and weakening generalisation. Separating the representations makes the actor representation more instance-agnostic and improves both transfer and learning efficiency by reducing interference between the actor and critic objectives. It also proves more parameter-efficient than increasing the capacity of a shared representation.

The chapter also showed that architectural decoupling does not eliminate all interaction between the actor and critic. Applying the information metrics to the observations collected during training revealed that the critic influences the policy's state-visitation distribution even when the two networks share no representation parameters. Because the critic's value estimates enter the policy-gradient objective, they shape policy updates and, consequently, the experience collected by the actor. This provides evidence that the critic plays an implicit role in exploration, potentially steering the policy towards data that facilitates value learning. However, the same mechanism may introduce residual objective interference when the experience favoured by the critic does not support learning an optimal, transferable policy. The chapter therefore identifies that the critic actively influences the data collected by the agent.

\paragraph{Chapter~3: Scenario generation improves transfer only when it remains grounded in the target distribution.}
Chapter~3 moved from the representations learned by the agent to the distribution of environments used during training. It first applied the representational analysis developed in Chapter~2 to adaptive level sampling. Prioritising instances according to value loss was found to reduce the information that shared actor--critic representations encode about training-instance identity. This result provides a representational explanation for the regularising effect of value-loss-prioritised sampling: by reducing the sampling frequency of instances with low value loss, the strategy progressively deprioritises environments to which the agent's learned representation has already become overfitted.

The chapter then examined methods that generate new environments and demonstrated that novelty alone does not imply useful diversity. Existing unsupervised environment-design methods can produce a substantial shift between the distribution used for training and the distribution on which zero-shot transfer is evaluated. Optimising generated environments primarily according to their difficulty for the current agent may increasingly concentrate training on unusual, degenerate, or adversarial scenarios. Although these scenarios create learning pressure, they need not teach behaviours relevant to the target task and can ultimately reduce transfer.

To address this failure, the chapter introduced Data-Regularised Environment Design (\hbox{DRED}). DRED learns an approximation to the target instance distribution from an initially available dataset of environment parameters and samples generated instances from this learned distribution through an adaptive training procedure. The resulting environments extend beyond the finite dataset while retaining its semantic and distributional structure. Across the evaluated domains, DRED promotes learning instance-invariant representations and improves zero-shot generalisation relative to fixed-set sampling, adaptive level-replay methods, and prior unsupervised environment-design approaches. The principal finding is therefore that effective environment generation must balance two requirements: it must expand the training distribution, but remain grounded in the distribution to which the policy is expected to transfer.

\paragraph{Chapter~4: Persistent 3D state improves the coherence and controllability of learned simulators.}
Chapter~4 addressed the first and more fundamental condition of the simulation bottleneck: access to a simulator of the target environment. It introduced \hbox{PERSIST}, a framework that learns a simulator directly from offline interaction data. PERSIST generates action-conditioned visual experience by maintaining and evolving a persistent latent 3D state. Existing interactive video models commonly encode the environment implicitly through a bounded history of pixel observations. Once observations leave this temporal context, information about occluded objects, previously visited locations, or earlier modifications to the environment may become inaccessible. As a result, geometry can drift, objects can disappear, and revisited regions can become inconsistent with the history of interaction.

PERSIST supplements bounded pixel histories with a persistent 3D representation of the environment. The 3D state is updated as the agent interacts with the environment and is projected into screen space to condition visual generation. This allows spatial information to be retrieved according to the agent's current position and orientation rather than reconstructed repeatedly from a short sequence of recent frames. Quantitative evaluations and a qualitative user study show that this design improves spatial memory, geometric consistency, and long-horizon stability relative to approaches without persistent spatial state. In particular, PERSIST sustains coherent, evolving environments over interactions hundreds of times longer than the visual generator's immediate context window.

The explicit 3D state also introduces forms of control that are difficult to realise in models whose world representation remains implicit in pixels. PERSIST can synthesise diverse 3D environments from a single image and supports environment specification and editing through direct manipulation of the underlying spatial representation. These capabilities indicate that structured world models can provide more than improved prediction: their internal state can serve as an interface for constructing, inspecting, and controlling generated environments. PERSIST therefore advances the feasibility of neural simulators learned from offline interaction data, while demonstrating that persistent spatial state is an important inductive bias for models intended to support extended interaction.

\paragraph{Reducing the simulation bottleneck requires controlling information and distributions at multiple levels.}
Taken together, the prior chapters establish a hierarchy of interventions on the simulation bottleneck. Chapter~2 acts at the level of the agent, showing that architectural and representational choices determine whether experience is compressed into transferable features or into features tied to particular training environments. Chapter~3 acts at the level of the training distribution, showing that generative augmentation can compensate for limited scenario coverage when generation is constrained by deployment-relevant data. Chapter~4 acts at the level of the simulator, showing that learned action-conditioned experience benefits from a persistent state structured around the spatial organisation of the environment.


\section{Limitations and Directions for Future Research}\label{sec:limit}

% 3 - 5 subsections with 1 (sometimes 2) paragraphs outlining limitations and directions of future work. 

% 1. Neural-symbolic simulator hybrids
% plan:
% - DRED generates simulator input parameters but relies on a traditional simulator to generate interaction data.
% - In contrast PERSIST is a entirely replaces the similator with a fully neural pipeline that simulates environment dynamics and produces a 3D state and RGB observations.
% - While it greatly improves stability and consistency PERSIST only support minutes-long generation and does not have perfect consistency, as traditional simulators do.
% - One can imagine an hybrid version, for example using a neural component for photo-realistic rendering but rely on a traditional simulator to predict complex dynamics such as contact forces \cite{Mujoco}. Recent projects from industry labs seem to be pointing towards that direction \cite{Roblox, project EDEN}
\subsection{Hybrid neural and structured simulators}

DRED and PERSIST intervene at opposite ends of the simulation pipeline. DRED generates parameters for new environment instances but relies on a conventional simulator to execute them and produce interaction data. PERSIST instead replaces the simulator with a fully neural pipeline that predicts environment dynamics, maintains a latent 3D state, and generates RGB observations. Although this approach substantially improves the consistency and duration of neural simulation, PERSIST remains less reliable than a conventional simulator. Its generated worlds remain coherent for minutes rather than indefinitely, and its state can still accumulate errors or depart from the rules of the underlying environment. Conventional simulators, by contrast, explicitly enforce their state-transition rules, but require those rules and the corresponding assets to be specified manually.

A promising direction is therefore to combine neural and structured simulation rather than treating them as mutually exclusive alternatives. A conventional simulator could maintain quantities for which precision is essential, such as object identities, collision geometry, contact forces, and other physical constraints \citep{TodorovET12}, while neural components model aspects that are difficult or expensive to specify, such as photorealistic appearance, sensor observations, or residual dynamics. Such a system would retain the long-horizon consistency and controllability of an explicit simulator while reducing the burden of constructing high-fidelity visual environments. Recent hybrid systems combine physics-based simulation with generative rendering, or maintain an authoritative structured world state while using video models to produce visual observations \citep{mo2025dreamland,singh2026robloxreality}. Related projects separate persistent scene state from visual generation while retaining a learned representation of the world \citep{vastai2026projecteden}. Developing principled interfaces between these components---including how neural predictions should update structured state and how disagreements between them should be resolved---is an important direction for learned simulation.

% 2. Extensions to PERSIST [needs different title, maybe extending neural simulator capabilities?]
% - 3D memory bank: load/unload spatial chunks from gpu memory, exploiting the 3D structure of the latent representation
% - Multiplayer \cite{MultiGen}: exploit factorization of dynamics/renderer to have a single dynamics module attending to many (1 per player) renderers.
\subsection{Scaling neural simulation to larger worlds and multiple agents}

PERSIST maintains spatial memory within a fixed-size 3D volume centred on the agent. This bounds its computational and memory requirements, but causes information to be discarded as regions move outside the active volume, preventing full spatial consistency across arbitrarily large worlds. A natural extension would be to maintain a spatial memory bank composed of smaller spatial chunks. Chunks close to the agent, visible from its current viewpoint, or likely to be affected by an action could be loaded into accelerator memory, while remaining chunks exist in lower-cost storage. When the agent returns to a previously visited location, the corresponding spatial chunk could be retrieved and integrated into the active representation. Hierarchical spatial resolution could further allocate more capacity to nearby or interaction-relevant regions while retaining only coarse information elsewhere. The main challenge would be to preserve consistency across chunk boundaries and to model meaningful changes of chunks that remain unobserved for long periods. Solving this problem would extend persistent neural simulation over arbitrarily large worlds.

The factorisation of environment state, camera dynamics, and rendering also provides a natural foundation for multi-agent simulation. Rather than maintaining an independent world for each player, a multiplayer extension could evolve one shared environment state using the joint actions of all participants, while associating each participant with a separate camera and rendering process. Every player would then receive an egocentric observation conditioned on the same underlying world state, allowing actions performed from one viewpoint to remain visible and consequential from the others. Recent work on multiplayer world models similarly uses shared persistent memory to coordinate several viewpoints \citep{po2026multigen}. Extending PERSIST in this direction would require models that resolve simultaneous and potentially conflicting actions, preserve agent identities, and maintain consistency across observations generated for different players. Beyond games, these capabilities would support learned simulation of multi-robot systems and other environments in which several agents interact within the same environment.

% 3. World Action Models
% - Acknowledge PERSIST falls short of training an RL agent. 
% - Naive approach of employing neural simulators: data generation
% - More exciting: distill agentic capabilities into a neural simulator. Allows agent to bootstrap from a structured, compact representation with learned priors. (refer to DEAC[chapter 2] on how important learned representations are to obtain generalisable policies). Concretely, this is as easy as incorporating an action prediction head into the model, and provide task instructions as part of the prompt. At deployment, the action prediction head is used as a policy. The model conditions on real observation feedback from the environment, its own (state/obs) predictions are discarded.
% - [Note : mention VLAs? Mention use for planning at deployment?]
% - Draw parallels with LLM pre-training (learn simulator of language) and Post-Training: give LLM agentic capabilities to manipulate their language outputs to achieve goals, e.g. becoming a chatbot, coding agent. 
% Expect WAM to follow a similar path to LLMs, e.g. SFT\cite{GPT-instruct/3.5} -> RLHF & RLVR . \cite{Dreamer-v1,v4} -> training in-imagination works and benefits from scale.
% WAMs can already achieve short-horizon tasks \cite{Cosmos-PolicyPaper}. [Note: finding the VLA vs WAM study would be helpful here]. Achieving long-horizon neural simulation opens up learning new long-horizon tasks for the resulting agents.
% [note: maybe needs to be more concrete/narrow focused in future directions? What specifically shall we do to give agentic capabilities to neural sims?]
\subsection{From world models to world action models}
The most important limitation of PERSIST is that it does not yet demonstrate the training of an agent within its generated experience. The direct use of a neural simulator would be to generate trajectories for planning or reinforcement learning, replacing interaction with either a hand-engineered simulator or the physical environment. Prior model-based RL methods have shown that policies can be learned through imagined experience \citep{hafner2025mastering,alonso2024diffusion,micheli2023transformers}, and recent work has demonstrated that this approach can scale to complex domains such as Minecraft \citep{hafner2025training}. Applying this paradigm to PERSIST would require more than visually coherent generation. The model would need to preserve task-relevant action consequences accurately enough that policies cannot exploit simulation errors, and its inference cost would need to permit the large number of online interactions required for RL.

A more direct direction is to post-train the world model itself into a policy, producing a \emph{world action model}. This parallels the development of large language models. Predictive pre-training first produces a model of the distribution and dynamics of language, after which supervised and reinforcement-learning post-training convert that model into a system capable of pursuing user-specified goals \citep{ouyang2022training}. Analogously, world-model pre-training can provide a broad model of visual dynamics, while subsequent action supervision and reinforcement learning can teach the model to manipulate those dynamics in pursuit of a task. Recent world action models have already demonstrated that video-model representations can be adapted for short-horizon visuomotor control and model-based planning \citep{kim2026cosmos}, with early comparisons suggesting that their predictive pre-training can achieve stronger generalisation than vision--language--action models \citep{wang2026world}. %todo:check claim 
Extending persistent world models to longer-horizon control would test whether the advantages of a persistent state carry over from simulation quality to policy learning. In particular, it would establish whether agents can exploit the model's explicit spatial memory to complete tasks that require returning to earlier locations, manipulating objects over extended sequences, or reasoning about portions of the environment outside the current field of view.